\documentclass[11pt]{article}
\usepackage{amsmath, amssymb, amsthm, algorithm, algorithmic}
\usepackage[margin=1in]{geometry}
\usepackage{xcolor}

\title{Convergence Analysis of Stochastic Variance Reduced Gradient (SVRG)}
\date{}

\begin{document}
\maketitle

\section*{Algorithm Name}
Stochastic Variance Reduced Gradient (SVRG)

\section*{Mathematical Formulation}
The SVRG algorithm is designed to minimize a finite-sum objective function:
\begin{equation}
    \min_{x \in \mathbb{R}^d} f(x) = \frac{1}{n} \sum_{i=1}^n f_i(x)
\end{equation}
To achieve this, SVRG utilizes a double-loop structure to periodically compute a full gradient and use it to reduce the variance of the stochastic gradients in the inner loop. 

Let $\tilde{x}_{s-1}$ be the anchor point at the outer loop stage $s$. The algorithm computes the exact full gradient at this anchor point:
\begin{equation}
    \tilde{\mu}_s = \nabla f(\tilde{x}_{s-1}) = \frac{1}{n} \sum_{i=1}^n \nabla f_i(\tilde{x}_{s-1})
\end{equation}
In the inner loop, for iterations $t = 1, \dots, m$, a random index $i_t$ is sampled uniformly from $\{1, \dots, n\}$. The variance-reduced stochastic gradient estimator is constructed as:
\begin{equation}
    v_t = \nabla f_{i_t}(x_{t-1}) - \nabla f_{i_t}(\tilde{x}_{s-1}) + \tilde{\mu}_s
\end{equation}
The parameter update is then performed using a constant learning rate $\eta$:
\begin{equation}
    x_t = x_{t-1} - \eta v_t
\end{equation}

\section*{Pseudocode}
\begin{algorithm}[H]
\caption{Stochastic Variance Reduced Gradient (SVRG)}
\begin{algorithmic}[1]
\REQUIRE Learning rate $\eta > 0$, inner loop length $m$, number of stages $S$, initial point $\tilde{x}_0$
\FOR{$s = 1, 2, \dots, S$}
    \STATE Compute the full gradient at the anchor point: $\tilde{\mu}_s = \frac{1}{n} \sum_{i=1}^n \nabla f_i(\tilde{x}_{s-1})$
    \STATE Initialize the inner loop: $x_0 = \tilde{x}_{s-1}$
    \FOR{$t = 1, 2, \dots, m$}
        \STATE Sample index $i_t \sim \text{Uniform}(\{1, \dots, n\})$
        \STATE Compute variance-reduced gradient: $v_t = \nabla f_{i_t}(x_{t-1}) - \nabla f_{i_t}(\tilde{x}_{s-1}) + \tilde{\mu}_s$
        \STATE Update iterate: $x_t = x_{t-1} - \eta v_t$
    \ENDFOR
    \STATE Set new anchor point $\tilde{x}_s$ by choosing uniformly at random from $\{x_0, x_1, \dots, x_{m-1}\}$
\ENDFOR
\RETURN $\tilde{x}_S$
\end{algorithmic}
\end{algorithm}

\section*{Dry Run Example}
Let us apply SVRG to $f(w) = (w-3)^2$ with $w_0 = 0$, $\eta = 0.1$, and $m = 3$ inner iterations. Because $f(w)$ is a single function, we treat it as $n=1$ where $f_1(w) = (w-3)^2$. Note that for $n=1$, SVRG reduces exactly to standard Gradient Descent because the stochastic variance trivially cancels out.

\textbf{Initialization:} 
Anchor $\tilde{w}_0 = 0$. Full gradient $\tilde{\mu}_1 = \nabla f(\tilde{w}_0) = 2(0-3) = -6$. Let $w_0 = \tilde{w}_0 = 0$.

\textbf{Inner Loop $t=1$:}
\begin{itemize}
    \item Sample $i_1 = 1$.
    \item Stochastic gradients: $\nabla f_1(w_0) = 2(0-3) = -6$, $\nabla f_1(\tilde{w}_0) = -6$.
    \item Variance-reduced gradient: $v_1 = -6 - (-6) + (-6) = -6$.
    \item Update: $w_1 = w_0 - \eta v_1 = 0 - 0.1(-6) = 0.6$.
    \item Objective value: $f(w_1) = (0.6 - 3)^2 = 5.76$.
\end{itemize}

\textbf{Inner Loop $t=2$:}
\begin{itemize}
    \item Sample $i_2 = 1$.
    \item Stochastic gradients: $\nabla f_1(w_1) = 2(0.6-3) = -4.8$.
    \item Variance-reduced gradient: $v_2 = -4.8 - (-6) + (-6) = -4.8$.
    \item Update: $w_2 = w_1 - \eta v_2 = 0.6 - 0.1(-4.8) = 1.08$.
    \item Objective value: $f(w_2) = (1.08 - 3)^2 = 3.6864$.
\end{itemize}

\textbf{Inner Loop $t=3$:}
\begin{itemize}
    \item Sample $i_3 = 1$.
    \item Stochastic gradients: $\nabla f_1(w_2) = 2(1.08-3) = -3.84$.
    \item Variance-reduced gradient: $v_3 = -3.84 - (-6) + (-6) = -3.84$.
    \item Update: $w_3 = w_2 - \eta v_3 = 1.08 - 0.1(-3.84) = 1.464$.
    \item Objective value: $f(w_3) = (1.464 - 3)^2 \approx 2.359$.
\end{itemize}

\section*{Assumptions}
The convergence analysis relies on the following standard assumptions:
\begin{enumerate}
    \item \textbf{$L$-Smoothness:} Each component function $f_i(x)$ is continuously differentiable and $L$-smooth. For all $x, y \in \mathbb{R}^d$:
    \begin{equation}
        \|\nabla f_i(x) - \nabla f_i(y)\| \leq L \|x - y\|
    \end{equation}
    \item \textbf{$\mu$-Strong Convexity:} The full objective function $f(x)$ is $\mu$-strongly convex. For all $x, y \in \mathbb{R}^d$:
    \begin{equation}
        f(y) \geq f(x) + \langle \nabla f(x), y - x \rangle + \frac{\mu}{2} \|y - x\|^2
    \end{equation}
    \item \textbf{Optimal Solution:} There exists a unique optimal solution $x^*$ such that $\nabla f(x^*) = 0$.
\end{enumerate}

\section*{Main Convergence Theorem}
\textbf{Theorem (Linear Convergence of SVRG):} Suppose the assumptions of $L$-smoothness and $\mu$-strong convexity hold. If the learning rate $\eta$ and the inner loop length $m$ are chosen such that $0 < \eta < \frac{1}{2L}$ and 
\begin{equation}
    \alpha = \frac{1}{\mu \eta (1 - 2L\eta) m} + \frac{2L\eta}{1 - 2L\eta} < 1
\end{equation}
then the SVRG algorithm converges linearly in expectation:
\begin{equation}
    \mathbb{E}[f(\tilde{x}_s) - f(x^*)] \leq \alpha^s [f(\tilde{x}_0) - f(x^*)]
\end{equation}

\section*{Theoretical Properties}
\begin{itemize}
    \item \textbf{Variance Reduction:} Unlike traditional SGD where the variance of the gradient estimator prevents convergence with a constant learning rate, the SVRG estimator $v_t$ satisfies $\mathbb{E}[\|v_t\|^2] \to 0$ as $x_t \to x^*$ and $\tilde{x} \to x^*$.
    \item \textbf{Constant Learning Rate:} Due to the variance reduction property, SVRG achieves linear convergence to the exact optimum using a fixed step size, eliminating the need for delicate learning rate decay schedules.
    \item \textbf{Gradient Complexity:} SVRG requires $O((n + \frac{L}{\mu}) \log(1/\epsilon))$ component gradient evaluations to reach $\epsilon$-accuracy, which is significantly faster than the $O(n \frac{L}{\mu} \log(1/\epsilon))$ of full Gradient Descent.
\end{itemize}

\section*{Discussion}
SVRG represents a major breakthrough in first-order optimization for finite sums. By storing the full gradient at an anchor point, it corrects the trajectory of the inner stochastic updates. The double-loop structure provides a trade-off: the outer loop evaluates $n$ gradients, which is expensive, but amortizes this cost over $m$ cheap inner iterations. Setting $m = O(n)$ effectively balances the computational load, yielding an algorithm that enjoys the cheap per-iteration cost of SGD but the robust linear convergence rate of full Gradient Descent.

\section*{Preliminaries (Definitions and Lemmas)}

Before the main proof, we establish two critical lemmas.

\textbf{Lemma 1 (Co-coercivity of Smooth Functions).} 
For an $L$-smooth function $f_i$ optimized at $x^*$, the following bound holds:
\begin{equation}
    \mathbb{E}_{i}[\|\nabla f_i(x) - \nabla f_i(x^*)\|^2] \leq 2L (f(x) - f(x^*))
\end{equation}
\textit{Proof of Lemma 1:}
\begin{itemize}
    \item \textbf{[Step 1]} Define the shifted function $g_i(x) = f_i(x) - f_i(x^*) - \langle \nabla f_i(x^*), x - x^* \rangle$. Note that $\nabla g_i(x) = \nabla f_i(x) - \nabla f_i(x^*)$.
    \item \textbf{[Step 2]} Since $f_i$ is $L$-smooth, $g_i$ is also $L$-smooth. By the standard descent lemma, for any $x, y$: $g_i(y) \leq g_i(x) + \langle \nabla g_i(x), y - x \rangle + \frac{L}{2} \|y - x\|^2$.
    \item \textbf{[Step 3]} Minimize both sides with respect to $y$. The minimum of the right hand side occurs at $y = x - \frac{1}{L} \nabla g_i(x)$.
    \item \textbf{[Step 4]} Substitute $y = x - \frac{1}{L} \nabla g_i(x)$ into the descent inequality to get: $g_i(y) \leq g_i(x) - \frac{1}{L} \|\nabla g_i(x)\|^2 + \frac{L}{2} \|-\frac{1}{L} \nabla g_i(x)\|^2$.
    \item \textbf{[Step 5]} Simplify the right side: $g_i(y) \leq g_i(x) - \frac{1}{2L} \|\nabla g_i(x)\|^2$.
    \item \textbf{[Step 6]} Since $g_i$ is convex and its minimum is 0 (at $x^*$), $g_i(y) \geq 0$. Thus, $\frac{1}{2L} \|\nabla g_i(x)\|^2 \leq g_i(x)$.
    \item \textbf{[Step 7]} Substitute back the definitions of $g_i(x)$ and $\nabla g_i(x)$: $\frac{1}{2L} \|\nabla f_i(x) - \nabla f_i(x^*)\|^2 \leq f_i(x) - f_i(x^*) - \langle \nabla f_i(x^*), x - x^* \rangle$.
    \item \textbf{[Step 8]} Take the expectation over the uniform random choice of index $i \in \{1, \dots, n\}$. Since $\mathbb{E}_i[f_i(x)] = f(x)$, we have: $\frac{1}{2L} \mathbb{E}_i[\|\nabla f_i(x) - \nabla f_i(x^*)\|^2] \leq f(x) - f(x^*) - \langle \nabla f(x^*), x - x^* \rangle$.
    \item \textbf{[Step 9]} By the optimality of $x^*$, the full gradient $\nabla f(x^*) = 0$. Thus, we obtain $\mathbb{E}_i[\|\nabla f_i(x) - \nabla f_i(x^*)\|^2] \leq 2L(f(x) - f(x^*))$.
\end{itemize}

\textbf{Lemma 2 (Variance Bound of SVRG Estimator).}
The variance-reduced gradient $v_t$ satisfies:
\begin{equation}
    \mathbb{E}_{i_t}[\|v_t\|^2] \leq 4L(f(x_{t-1}) - f(x^*)) + 4L(f(\tilde{x}) - f(x^*))
\end{equation}
\textit{Proof of Lemma 2:}
\begin{itemize}
    \item \textbf{[Step 10]} Expand the definition of the estimator $v_t$: $\mathbb{E}[\|v_t\|^2] = \mathbb{E}[\|\nabla f_{i_t}(x_{t-1}) - \nabla f_{i_t}(\tilde{x}) + \nabla f(\tilde{x})\|^2]$.
    \item \textbf{[Step 11]} Add and subtract the optimal gradient $\nabla f_{i_t}(x^*)$ inside the norm: $\mathbb{E}[\|v_t\|^2] = \mathbb{E}[\|(\nabla f_{i_t}(x_{t-1}) - \nabla f_{i_t}(x^*)) - (\nabla f_{i_t}(\tilde{x}) - \nabla f_{i_t}(x^*) - \nabla f(\tilde{x}))\|^2]$.
    \item \textbf{[Step 12]} Apply the inequality $\|a - b\|^2 \leq 2\|a\|^2 + 2\|b\|^2$: $\mathbb{E}[\|v_t\|^2] \leq 2\mathbb{E}[\|\nabla f_{i_t}(x_{t-1}) - \nabla f_{i_t}(x^*)\|^2] + 2\mathbb{E}[\|\nabla f_{i_t}(\tilde{x}) - \nabla f_{i_t}(x^*) - \nabla f(\tilde{x})\|^2]$.
    \item \textbf{[Step 13]} Recognize that $\nabla f(x^*) = 0$, so $\nabla f(\tilde{x}) = \nabla f(\tilde{x}) - \nabla f(x^*) = \mathbb{E}[\nabla f_{i_t}(\tilde{x}) - \nabla f_{i_t}(x^*)]$. The second term is thus the variance of $\nabla f_{i_t}(\tilde{x}) - \nabla f_{i_t}(x^*)$.
    \item \textbf{[Step 14]} Use the standard variance bound $\mathbb{E}[\|X - \mathbb{E}[X]\|^2] \leq \mathbb{E}[\|X\|^2]$ on the second term to get: $\mathbb{E}[\|\nabla f_{i_t}(\tilde{x}) - \nabla f_{i_t}(x^*) - \nabla f(\tilde{x})\|^2] \leq \mathbb{E}[\|\nabla f_{i_t}(\tilde{x}) - \nabla f_{i_t}(x^*)\|^2]$.
    \item \textbf{[Step 15]} Combine the bounds to obtain: $\mathbb{E}[\|v_t\|^2] \leq 2\mathbb{E}[\|\nabla f_{i_t}(x_{t-1}) - \nabla f_{i_t}(x^*)\|^2] + 2\mathbb{E}[\|\nabla f_{i_t}(\tilde{x}) - \nabla f_{i_t}(x^*)\|^2]$.
    \item \textbf{[Step 16]} Apply Lemma 1 to both terms, yielding the final variance bound: $\mathbb{E}[\|v_t\|^2] \leq 2(2L(f(x_{t-1}) - f(x^*))) + 2(2L(f(\tilde{x}) - f(x^*))) = 4L(f(x_{t-1}) - f(x^*)) + 4L(f(\tilde{x}) - f(x^*))$.
\end{itemize}

\section*{Main Proof (Step-by-Step Derivation)}

We analyze a single inner loop (stage $s$). Let $\tilde{x} = \tilde{x}_{s-1}$ be the anchor point.
\begin{itemize}
    \item \textbf{[Step 17]} Start with the squared Euclidean distance to the optimum at iteration $t$: $\|x_t - x^*\|^2 = \|x_{t-1} - \eta v_t - x^*\|^2$.
    \item \textbf{[Step 18]} Expand the squared norm: $\|x_t - x^*\|^2 = \|x_{t-1} - x^*\|^2 - 2\eta \langle v_t, x_{t-1} - x^* \rangle + \eta^2 \|v_t\|^2$.
    \item \textbf{[Step 19]} Take the expectation conditioned on $x_{t-1}$ with respect to the random choice of $i_t$: $\mathbb{E}_{i_t}[\|x_t - x^*\|^2] = \|x_{t-1} - x^*\|^2 - 2\eta \langle \mathbb{E}_{i_t}[v_t], x_{t-1} - x^* \rangle + \eta^2 \mathbb{E}_{i_t}[\|v_t\|^2]$.
    \item \textbf{[Step 20]} Since $\mathbb{E}_{i_t}[\nabla f_{i_t}(x_{t-1}) - \nabla f_{i_t}(\tilde{x})] = \nabla f(x_{t-1}) - \nabla f(\tilde{x})$, the estimator is unbiased: $\mathbb{E}_{i_t}[v_t] = \nabla f(x_{t-1})$. Substitute this into the inner product: $\mathbb{E}_{i_t}[\|x_t - x^*\|^2] = \|x_{t-1} - x^*\|^2 - 2\eta \langle \nabla f(x_{t-1}), x_{t-1} - x^* \rangle + \eta^2 \mathbb{E}_{i_t}[\|v_t\|^2]$.
    \item \textbf{[Step 21]} By the convexity of $f$, we have $f(x^*) \geq f(x_{t-1}) + \langle \nabla f(x_{t-1}), x^* - x_{t-1} \rangle$, which rearranges to $-\langle \nabla f(x_{t-1}), x_{t-1} - x^* \rangle \leq -(f(x_{t-1}) - f(x^*))$.
    \item \textbf{[Step 22]} Substitute the convexity bound into the distance expansion: $\mathbb{E}_{i_t}[\|x_t - x^*\|^2] \leq \|x_{t-1} - x^*\|^2 - 2\eta (f(x_{t-1}) - f(x^*)) + \eta^2 \mathbb{E}_{i_t}[\|v_t\|^2]$.
    \item \textbf{[Step 23]} Substitute the variance bound from Lemma 2: $\mathbb{E}_{i_t}[\|x_t - x^*\|^2] \leq \|x_{t-1} - x^*\|^2 - 2\eta (f(x_{t-1}) - f(x^*)) + 4L\eta^2 (f(x_{t-1}) - f(x^*)) + 4L\eta^2 (f(\tilde{x}) - f(x^*))$.
    \item \textbf{[Step 24]} Group the terms involving $(f(x_{t-1}) - f(x^*))$: $\mathbb{E}_{i_t}[\|x_t - x^*\|^2] \leq \|x_{t-1} - x^*\|^2 - 2\eta(1 - 2L\eta) (f(x_{t-1}) - f(x^*)) + 4L\eta^2 (f(\tilde{x}) - f(x^*))$.
    \item \textbf{[Step 25]} Take the total expectation and sum the inequality over the inner loop from $t=1$ to $m$: $\sum_{t=1}^m \mathbb{E}[\|x_t - x^*\|^2] \leq \sum_{t=1}^m \mathbb{E}[\|x_{t-1} - x^*\|^2] - 2\eta(1 - 2L\eta) \sum_{t=1}^m \mathbb{E}[f(x_{t-1}) - f(x^*)] + 4Lm\eta^2 (f(\tilde{x}) - f(x^*))$.
    \item \textbf{[Step 26]} Telescope the sum $\sum_{t=1}^m (\mathbb{E}[\|x_t - x^*\|^2] - \mathbb{E}[\|x_{t-1} - x^*\|^2]) = \mathbb{E}[\|x_m - x^*\|^2] - \mathbb{E}[\|x_0 - x^*\|^2]$, which yields: $\mathbb{E}[\|x_m - x^*\|^2] \leq \mathbb{E}[\|x_0 - x^*\|^2] - 2\eta(1 - 2L\eta) \sum_{t=1}^m \mathbb{E}[f(x_{t-1}) - f(x^*)] + 4Lm\eta^2 (f(\tilde{x}) - f(x^*))$.
    \item \textbf{[Step 27]} Since the squared distance $\mathbb{E}[\|x_m - x^*\|^2] \geq 0$, we can drop it to form a strict inequality and rearrange the terms: $2\eta(1 - 2L\eta) \sum_{t=1}^m \mathbb{E}[f(x_{t-1}) - f(x^*)] \leq \mathbb{E}[\|x_0 - x^*\|^2] + 4Lm\eta^2 (f(\tilde{x}) - f(x^*))$.
    \item \textbf{[Step 28]} Divide the entire inequality by $2\eta(1 - 2L\eta)m$: $\frac{1}{m} \sum_{t=1}^m \mathbb{E}[f(x_{t-1}) - f(x^*)] \leq \frac{\mathbb{E}[\|x_0 - x^*\|^2]}{2\eta(1 - 2L\eta)m} + \frac{2L\eta}{1 - 2L\eta} (f(\tilde{x}) - f(x^*))$.
    \item \textbf{[Step 29]} Because the new anchor point $\tilde{x}_s$ is chosen uniformly at random from $\{x_0, \dots, x_{m-1}\}$, its expected objective gap is exactly the left-hand side: $\mathbb{E}[f(\tilde{x}_s) - f(x^*)] = \frac{1}{m} \sum_{t=1}^m \mathbb{E}[f(x_{t-1}) - f(x^*)]$.
    \item \textbf{[Step 30]} By the $\mu$-strong convexity of $f$, we have $f(\tilde{x}) - f(x^*) \geq \frac{\mu}{2} \|\tilde{x} - x^*\|^2$. Since the inner loop is initialized with $x_0 = \tilde{x}$, we can bound the initial distance as $\mathbb{E}[\|x_0 - x^*\|^2] \leq \frac{2}{\mu} (f(\tilde{x}) - f(x^*))$.
    \item \textbf{[Step 31]} Substitute this strong convexity bound into the inequality from Step 28: $\mathbb{E}[f(\tilde{x}_s) - f(x^*)] \leq \left( \frac{2}{2\mu\eta(1 - 2L\eta)m} + \frac{2L\eta}{1 - 2L\eta} \right) (f(\tilde{x}) - f(x^*))$.
    \item \textbf{[Step 32]} Simplify the expression to arrive at the final linear convergence relationship: $\mathbb{E}[f(\tilde{x}_s) - f(x^*)] \leq \left( \frac{1}{\mu\eta(1 - 2L\eta)m} + \frac{2L\eta}{1 - 2L\eta} \right) (f(\tilde{x}_{s-1}) - f(x^*))$.
\end{itemize}

\section*{Rate Derivation (With Constants)}

The derived contraction factor is:
\begin{equation}
    \alpha = \frac{1}{\mu\eta(1 - 2L\eta)m} + \frac{2L\eta}{1 - 2L\eta}
\end{equation}
To ensure convergence, we must choose $\eta$ and $m$ such that $\alpha < 1$. 
Let us set the learning rate to $\eta = \frac{1}{10L}$. It follows that $1 - 2L\eta = 1 - 0.2 = 0.8$. Substituting this into $\alpha$:
\begin{equation}
    \alpha = \frac{1}{\mu \left(\frac{1}{10L}\right) (0.8) m} + \frac{0.2}{0.8} = \frac{12.5 L}{\mu m} + 0.25
\end{equation}
To achieve a contraction factor of $\alpha \leq 0.5$, we require:
\begin{equation}
    \frac{12.5 L}{\mu m} \leq 0.25 \implies m \geq 50 \frac{L}{\mu}
\end{equation}
Let $\kappa = \frac{L}{\mu}$ be the condition number. With $\eta = \frac{1}{10L}$ and $m = 50\kappa$, SVRG halves the objective gap in expectation every outer stage $s$.

To reach an $\epsilon$-accurate solution, SVRG requires $S = \log_2(1/\epsilon)$ stages. Each stage requires:
\begin{itemize}
    \item $n$ gradient evaluations to compute the full gradient $\tilde{\mu}_s$.
    \item $2m$ gradient evaluations in the inner loop (one for $\nabla f_{i_t}(x_{t-1})$ and one for $\nabla f_{i_t}(\tilde{x})$, though the latter can be cached).
\end{itemize}
Thus, the total number of gradient evaluations per stage is $n + 2m = n + 100\kappa$. The overall gradient complexity to reach $\epsilon$-accuracy is therefore:
\begin{equation}
    O\left((n + \kappa) \log\left(\frac{1}{\epsilon}\right)\right)
\end{equation}

\section*{Special Cases}
\begin{itemize}
    \item \textbf{Single Component ($n=1$):} When the sum consists of only a single function, SVRG reduces exactly to standard Gradient Descent. The stochastic gradient is the full gradient, and the variance reduction term perfectly cancels out: $v_t = \nabla f(x_{t-1}) - \nabla f(\tilde{x}) + \nabla f(\tilde{x}) = \nabla f(x_{t-1})$.
    \item \textbf{Single Inner Step ($m=1$):} If the inner loop is run for only one iteration, the anchor point is updated immediately. Since $x_0 = \tilde{x}$, the first update is $v_1 = \nabla f_{i_1}(\tilde{x}) - \nabla f_{i_1}(\tilde{x}) + \tilde{\mu}_1 = \tilde{\mu}_1$. This is completely equivalent to full Gradient Descent.
    \item \textbf{Quadratic Functions:} For $f_i(x) = \frac{1}{2} x^T A_i x - b_i^T x$, the gradients are linear: $\nabla f_i(x) = A_i x - b_i$. The variance reduction term becomes $v_t = A_{i_t}(x_{t-1} - \tilde{x}_{s-1}) + \tilde{\mu}_s$. This avoids the need to evaluate the vector $b_{i_t}$ inside the inner loop, further saving computational resources.
\end{itemize}

\end{document}